\section{Related Work}

\paragraph{Target-guided and proactive dialogue planning.}
Target-guided dialogue studies how an agent can steer a conversation toward a
specified topic, target, or endpoint. Early methods model this process through
keyword prediction, discourse constraints, commonsense knowledge, and global
planning over knowledge graphs
\cite{tang-etal-2019-target,qin2020dynamic,zhong2021keyword,kishinami-etal-2022-target,yang-etal-2022-topkg}.
More recent proactive planners broaden the planning signal with stochastic path
modeling, action-topic planning, LLM-based reasoning, reflection, and
consistency correction
\cite{wang-etal-2023-dialogue,wang2024target,zheng-etal-2024-thoughts,guo-etal-2024-pcqpr,zhang-etal-2025-enhancing-goal}.
These approaches are closely related in spirit because they treat dialogue as a
long-horizon control problem rather than a sequence of independent responses.
Their targets, however, are usually specified before planning. ALMOND considers
a complementary setting in which the macro-goal remains expressed in language,
but the short-term skill needed to pursue it changes with the evolving dialogue
state.

\paragraph{Recommendation and negotiation foundations.}
Conversational recommendation provides data and tasks for proactive target
pursuit. DuRecDial and DuRecDial 2.0 introduce multi-domain recommendation
dialogues
\cite{liu-etal-2020-towards-conversational,liu-etal-2021-durecdial}, while
target-driven recommendation methods plan action-topic paths or retrieve past
experience to anticipate future states
\cite{wang2022follow,dao-etal-2024-experience}. Negotiation dialogue research
studies a different but equally strategic setting, where agents learn to bargain
through end-to-end optimization, self-play, or decoupled strategy and language
generation \cite{lewis-etal-2017-deal,he-etal-2018-decoupling}. These domains
motivate the skill space used in ALMOND, including bargaining, concession,
rapport maintenance, clarification, and deal-seeking behavior. They also
motivate the evaluation protocol: a dialogue may reach an acceptable final
outcome while still violating a budget constraint, taking too many turns, or
failing to adapt when the counterparty's intent shifts.

\paragraph{Adaptive dialogue policy learning.}
A growing line of work separates strategic dialogue planning from surface-form
generation. PPDPP trains a plug-and-play policy planner with goal-oriented
feedback \cite{deng2024plug}, and TRIP improves planning in non-collaborative
dialogue through diversified user simulation \cite{zhang-etal-2024-strength}.
PADPP is the closest multi-objective predecessor: it learns a planner that can
adapt to different numeric preference vectors at inference time
\cite{dao-liao-2025-one}. This makes PADPP a strong foundation for studying
preference-adaptive dialogue, but it also clarifies the interface that remains
unmodeled. In ALMOND, the control signal is not an observed vector; it is a
short-term skill selected under a natural-language macro-goal and revised as
the conversation changes. This shift is motivated by fine-grained intention
understanding and intent-drift detection
\cite{liang-etal-2025-intentionframe,wang2025detecting}, which emphasize that
dialogue state is not only informational but also strategic and affective.

\paragraph{Dynamic preferences and constrained skill learning.}
Outside dialogue, dynamic-weight MORL studies policies whose objective weights
change over time \cite{abels2019dynamic}, while generalized MORL learns
policies over a preference space \cite{yang2019generalized}. Surveys of MORL
clarify the trade-offs among utility functions, Pareto coverage, and planning
assumptions \cite{hayes2022practical}. Successor features and GPI provide a
mechanism for reusing value knowledge across reward functions
\cite{barreto2017successor,barreto2018transfer}, and constrained policy
optimization motivates treating safety requirements as constraints rather than
only as shaped rewards \cite{achiam2017constrained}. ALMOND draws on this view
of reusable behavior under changing objectives, but moves the adaptation signal
closer to dialogue practice: short-term skills are inferred from a linguistic
macro-goal, recent interaction history, and drift evidence. The regret-based
low-level mechanism further biases training toward skills that remain
under-mastered, allowing more complex communicative behavior to be bootstrapped
without retraining a new policy for every objective.
